\scanpage{12}

\mspage{6)}
La difficulté, c'est qu'il n'est pas clair que le foncteur

\[
 \widetilde C\longrightarrow W^{-1}C
\]

précédent soit aussi \uline{fidèle} (ce qui impliquerait que c'est une équivalence de catégories \inserted{et $=$ un iso}). \deleted{Dans ce cas, l'iso. inverse serait forcément donné par} \inserted{Il revient au $=$ de dire} que le foncteur can $C\longrightarrow\widetilde C$ (dont le composé avec $\widetilde C\longrightarrow W^{-1}C$ est bien le foncteur can. $C\longrightarrow W^{-1}C$) transforme flèches de $W$ en iso, ou encore que la construction qu'on vient \deleted{\unsure{de}} d'esquisser, pour \deleted{\unsure{le voir|les inv.}} \inserted{trouver} un inverse de $[f,i]$ dans $\widetilde C$, est bien valable, \deleted{\unsure{ce que}} Cela entraînerait d'ailleurs \inserted{semble-t-il,} que, \deleted{le} posant $W_0=W\cap\mathrm{cof}$, le foncteur

\[
 W_0^{-1}C\longrightarrow W^{-1}C
\]

serait une équivalence de catégories, car

\[
 \widetilde C\longrightarrow W^{-1}C
\]

se factorise visiblement par

\diagram{\begin{tikzcd}[ampersand replacement=\&,column sep=small,row sep=small]
\widetilde C \arrow[r] \arrow[d,"{\unsure{\wr|{\simeq}}}"{description},draw=none] \& W_0^{-1}C \arrow[r] \& W^{-1}C \\
W^{-1}C \& {} \& {}
\end{tikzcd}}

Mais je doute qu'il suffise d'inverser les équiv. faibles \uline{cofibrantes}, pour que ttes les équiv. faibles deviennent inversibles !

Je rectifie \deleted{donc} le tir, en prenant plutôt la catégorie (notée pourtant par $\widetilde C$) ayant $\mathord{=}^{\rm s}$ objets que $C$, \deleted{et où on désigne les flèches de $x$ à $y$} \inserted{et définie via les $\pi_0$ des $\underline{\mathrm{Ch}}_0(x,y)$}, \deleted{\unsure{définis|définies}} \inserted{cette catégorie} étant maintenant (pour $x,y$ donnés) \unsure{comme|connue} celle des diagrammes \deleted{\unsure{essayer}}

\diagram{\begin{tikzcd}[ampersand replacement=\&,column sep=small,row sep=small]
x \arrow[dr,tail,"f"'] \& y \arrow[d,"i"'," \sim"] \\
{} \& \bar y
\end{tikzcd}}

avec $\left\{\begin{aligned} f&\in\mathrm{cof}\\ i&\in W\end{aligned}\right.$
% archive sheet number: 11
